\chapter*{РЕФЕРАТ}
\addcontentsline{toc}{chapter}{РЕФЕРАТ}

Расчетно-пояснительная записка содержит 53 страницы, 12 рисунков, 1 таблицу, 15 источников.

КОМПЬЮТЕРНАЯ ГРАФИКА, ТРЁХМЕРНОЕ МОДЕЛИРОВАНИЕ, ВИЗУАЛИЗАЦИЯ, Z-БУФЕР, SHADOW MAPPING, ОСВЕЩЕНИЕ, ПОИСК ПУТИ, PYTHON, NUMPY.

Объект исследования --- процессы визуализации трёхмерных сцен и анимации движения объектов.

Цель работы --- разработка программного обеспечения для трёхмерного моделирования движения автомобиля в городской среде с использованием алгоритмов машинной графики.

В процессе работы был проведен анализ методов удаления невидимых поверхностей, алгоритмов закраски и построения теней. Обоснован выбор алгоритма построчного сканирования с Z-буфером и метода теневых карт. Спроектирована и реализована архитектура приложения на языке Python. Разработан графический конвейер, включающий матричные преобразования, растеризацию, текстурирование и расчет освещенности по модели Ламберта. Реализован модуль навигации на основе волнового алгоритма для ориентированного графа.

В результате создано интерактивное приложение, обеспечивающее корректную визуализацию городской среды с динамическими тенями и анимацией движения автомобиля. Исследование показало, что механизм отсечения невидимой геометрии демонстрирует высокую практическую значимость для оптимизации рендеринга. Анализ производительности выявил квадратичную зависимость времени рендеринга от коэффициента масштабирования сцены.